\documentclass{exam}
\usepackage{amsmath, amssymb}

\pagestyle{headandfoot}
\firstpageheadrule
\runningheadrule
\firstpageheader{Prof. Rong \\ Introduction to Real Analysis I}{Homework 4}{Aadhithya Saravanan}
\runningheader{Introduction to Real Analysis I \\ Homework 4}{}{Saravanan}
\firstpagefooter{}{}{}
\runningfooter{ }{\thepage}{ }

\printanswers

\begin{document}

\underline{Exercise 2.5.1}
\newline
\begin{parts}
    \part This is impossible. Let $(a_n)$ be the original sequence and $(a_{n_k})$ be the bounded subsequence. By the Bolzano-Weierstrass Theorem, the bounded subsequence $(a_{n_k})$ contains a subsequence that is convergent. This subsequence is also a subsequence of $(a_n)$ because each of its terms are in $(a_{n_k})$ and therefore also in $(a_n)$. So, this subsequence is a convergent subsequence of $(a_n)$, so it is not possible for a sequence to not have a convergent subsequence if it has a bounded subsequence.
    \part Let $a_n=\frac{1}{n}$ for odd $n$ and $a_n=1-\frac{1}{n}$ for even $n$. The subsequence defined with all odd $n$ converges to 0 and the subsequence defined with all even $n$ converges to 1.
    \part Let $(a_{1,n})$ be the sequence defined as $a_{1,n}=1-\frac{1}{n}$, converging to 1. Let $(a_{2,n})$ be the sequence defined as $a_{2,n}=\frac{1}{2}-\frac{1}{n}$, converging to $\frac{1}{2}$. In this way, we can define an infinite number of sequences, $(a_{k,n})$, each converging to $\frac{1}{k}$ for each $k\in \mathbb{N}$. Notice that the first $k$ terms of $(a_{k,n})$ are less than or equal to zero. We are going to define our main sequence $(x_n)$ as follows. The first term of $(x_n)$ is the first nonzero term of $(a_{1,n})$. Then, the second and third terms of $(x_n)$ are the second nonzero term of $(a_{1,n})$ and the first nonzero term of $(a_{2,n})$. We continue in this way, including the positive terms of $(a_{k,n})$ for all $k\in \mathbb{N}$ in $(x_n)$. Now, for any given element $\frac{1}{m}$ of $\{1, \frac{1}{2}, \frac{1}{3},\dots\}$, we can just look at the subsequence formed by $(a_{m,k})$ in $(x_n)$, which converges to $\frac{1}{m}$. Therefore, this $(x_n)$ contains subsequences that converge to every point of the infinite set $\{1, \frac{1}{2}, \frac{1}{3},\dots\}$.
    \part This is impossible. Any sequence with subsequences converging to every point of the infinite set $\{1, \frac{1}{2}, \frac{1}{3},\dots\}$ must also converge to 0, which is not in $\{1, \frac{1}{2}, \frac{1}{3},\dots\}$. By taking one element out of each of the subsequences that converge to $\frac{1}{n}$, we can construct a subsequence that approaches each element in $\{1, \frac{1}{2}, \frac{1}{3},\dots\}$. This subsequence converges to 0, as $\frac{1}{n}$ converges to 0. Therefore, it is not possible to have a sequence with subsequences converging to every element of $\{1, \frac{1}{2}, \frac{1}{3},\dots\}$ and no subsequence converging to a point outside this set.
    \newline
\end{parts}

\underline{Exercise 2.5.2}
\newline
\begin{parts}
    \part If every proper subsequence of $(x_n)$ converges, the subsequence that contains every element of $(x_n)$ except the first is convergent, as it is a proper subsequence. Let $a_n$ be this proper subsequence and it converges to $a$, so $a_n=x_{n+1}$. Then, by the definition of convergence, there exists an $N$ such that for $n\ge N$, $|a_n-a|<\epsilon$ for all $\epsilon>0$. Since $a_N=x_{N+1}$, we can let $N+1$ be our choice for which every term in $x_n$ for $n\ge N+1$ satisfies $|x_n-a|<\epsilon$. Since this condition is satisfied, $(x_n)$ converges and its limit is $a$.
    \part Let $(a_{n_k})$ be a divergent subsequence of $(x_n)$. That means there exists no $N$ such that $|a_{n_k}-a|<\epsilon$ for all $n_k\ge N$ and $a\in\mathbb{R}$. Since $(a_{n_k})$ has an infinite number of terms, there is no $M\in\mathbb{N}$ such that for all terms past $x_M$, there are no terms from $(a_{n_k})$ in the sequence. That means an infinite number of terms from $(a_{n_k})$ are past $x_M$ for any $M\in\mathbb{N}$. Since these terms from $(a_{n_k})$ are guaranteed to never be within a certain bound from a limit, the entire sequence past $x_M$ cannot be within a bound from a limit, so $(x_n)$ cannot be convergent.
    \part By the Bolzano-Weierstrass Theorem, since $(x_n)$ is bounded, there exists at least one convergent subsequence of $(x_n)$. Let this subsequence be $(x_{n_k})$ and let its limit be $L$. Let us assume that there does not exist two subsequences of $(x_n)$ that converge to different limits. This is the same as saying that all subsequence of $(x_n)$ converge to $L$. Since $(x_n)$ is divergent, we know it does not converge to any limit, including $L$. That means for some $\epsilon>0$, for all $N\in\mathbb{N}$, we can find an index $n\ge N$ such that $|x_n-L|\ge\epsilon$. Using this, we can construct a subsequence $(x_{m_j})$ such that $|x_{m_j}-L|\ge\epsilon$ for all $j$. This subsequence contains terms from $(x_n)$, so it is bounded. By the Bolzano-Weierstrass Theorem, it has a convergent subsequence $(x_{m_{j_l}})$ with a limit $M$. Every term in this subsequence satisfies $|x_{m_{j_l}}-L|\ge\epsilon$, so $|M-L|\ge\epsilon$ and $M\ne L$. But $(x_{m_{j_l}})$ is a subsequence of $(x_n)$ that does not converge to $L$. This contradicts that there does not exist two subsequences of $(x_n)$ that converge to different limits. Therefore, there must exist two subsequences of $(x_n)$ that converge to different limits.
    \part Let $(x_{n_k})$ be a convergent subsequence of $(x_n)$. Since there are an infinite number of terms in $(x_{n_k})$, we can always find a term in $(x_{n_k})$ that comes after any term in $(x_n)$. If that were not true, there would be a finite number of terms in $(x_{n_k})$. Since $(x_n)$ is monotone, we can say it is increasing without loss of generality. We have $x_n\le x_{n+1}$. Since there is always a term in $(x_{n_k})$ after any given $x_n$, we can choose $n_k\ge n$ and have $x_n\le x_{n_k}$. $(x_{n_k})$ is convergent, so all of its terms must be bounded, so $|x_{n_k}|\le M$ for some $M\in\mathbb{R}$. Then, we have $x_n\le x_{n_k}\le |x_{n_k}|\le M$. Then we can choose $K=\max\{|x_1|,|M|\}$ and we will have $|x_n|\le K$. So, $(x_n)$ is bounded. Since it is also monotone, $(x_n)$ is convergent.
    \newline
\end{parts}

\underline{Exercise 2.5.5}
\newline
\newline
We can prove this by contradiction. Assume that $(a_n)$ does not converge to $a$. This means that there exists an $\epsilon\ge0$ such that for all $N\in\mathbb{N}$, there exists an $n\ge N$ with $|a_n-a|\ge\epsilon$. The number of terms that satisfy this property must be infinite. If this were not true, then we could just choose a greater $N\in\mathbb{N}$ until there were no more terms outside of the range $(a-\epsilon, a+\epsilon)$. Since there are an infinite number of terms outside $(a-\epsilon, a+\epsilon)$, we can choose these terms to be a subsequence, which is also bounded because $(a_n)$ is bounded. Then, since this is a bounded sequence, it must contain a convergent subsequence by the Bolzano-Weierstrass Theorem.This subsequence cannot converge to $a$, as every term in the subsequence is outside of $(a-\epsilon, a+\epsilon)$. Since all of the terms of this subsequence are a part of the original $(a_n)$, this is a convergent subsequence of $(a_n)$ that does not converge to $a$. This is a contradiction, so our assumption that $(a_n)$ does not converge to $a$ is false. Therefore, $(a_n)$ must converge to $a$.
\newline

\underline{Exercise 2.5.9}
\newline
\newline
By definition of supremum, for every $\epsilon>0$, there exists an $x\in S$ such that $s-\frac{\epsilon}{2}<x$. Since $x\in S$, $x<a_n$ for infinitely many $a_n$. Then, $s-\frac{\epsilon}{2}<x<a_n$, which means any number less than $s$ is in $S$. By definition of supremum, we also know that $s+\frac{\epsilon}{2}\notin S$ for all $\epsilon>0$, so there are only finitely many $a_n$ satisfying $s+\frac{\epsilon}{2}<a_n$. Since there are an infinite number of $a_n$ satisfying $s-\frac{\epsilon}{2}<a_n$ and a finite number of $a_n$ satisfying $s+\frac{\epsilon}{2}<a_n$, there must be an infinite number of $a_n$ satisfying $s-\frac{\epsilon}{2}<a_n\le s+\frac{\epsilon}{2}$. So, for all $\epsilon>0$, there are infinitely many $a_n\in(s-\frac{\epsilon}{2}, s+\frac{\epsilon}{2}]$. Because $(s-\frac{\epsilon}{2},s+\frac{\epsilon}{2}]\subset(s-\epsilon, s+\epsilon)$, there are infinitely many $a_n\in(s-\epsilon, s+\epsilon)$. Let $\epsilon=\frac{1}{k}$. Then we can construct the subsequence which converges to $s$. As $k$ ranges over $\mathbb{N}$, let $a_{n_k}$ be an element of $(a_n)$ from the interval $(s-\frac{1}{k}, s+\frac{1}{k})$. This subsequence of $(a_n)$ converges to $s$, proving the Bolzano-Weierstrass Theorem using the Axiom of Completeness.
\newline

\underline{Exercise 2.6.2}
\newline
\begin{parts}
    \part A Cauchy sequence that is not monotone is $a_n=\frac{(-1)^n}{n+1}$. A sequence is Cauchy if for every $\epsilon>0$, there exists a $N\in\mathbb{N}$ such that whenever $m,n\ge N$, it follows that $|a_n-a_m|<\epsilon$. Choose $\epsilon=\frac{2}{N}$. Then, $|a_n-a_m|=|\frac{(-1)^n}{n+1}-\frac{(-1)^m}{m+1}|<|\frac{(-1)^n}{n}-\frac{(-1)^m}{m}|\le|\frac{1}{N}-\frac{-1}{N}|\le\frac{2}{N}$. So, $|a_n-a_m|<\frac{2}{N}=\epsilon$ and the sequence is Cauchy. Since its terms alternate signs, it cannot be monotone, so $(a_n)$ is a Cauchy sequence that is not monotone.
    \part If $(a_{n_k})$ is the unbounded subsequence of $(a_n)$, a Cauchy sequence, there does not exist a number $M>0$ such that $|a_{n_k}|<M$ for all $n_k\in\mathbb{N}$. Because all of the terms of $(a_{n_k})$ do not stay within a finite range, $(a_{n_k})$ cannot converge to a limit, so the subsequence is divergent. Since all of the terms of $(a_{n_k})$ are terms of $(a_n)$, that means an infinite number of terms of $(a_n)$ do not converge to a limit, so the entire sequence is divergent. By the Cauchy Criterion, $(a_n)$ does not converge, which is a contradiction. Therefore, there does not exist a Cauchy sequence with an unbounded subsequence.
    \part Assume there exists a divergent monotone sequence with a Cauchy subsequence. Let the sequence be $(a_n)$ and the Cauchy subsequence be $(a_{n_k})$ with limit $a$. Without loss of generality, let the sequence and subsequence be increasing. So, $a$ is an upper bound for $a_{n_k}$ for $n_k\in\mathbb{N}$. Since $(a_{n_k})$ is a subsequence of $(a_n)$, there must be an infinite number of terms of $(a_{n_k})$ in $(a_n)$. That means, for any $N\in\mathbb{N}$, there are always terms after $a_N$ that are a part of $(a_{n_k})$. So, since $(a_{n_k})$ is bounded above by $a$, $(a_n)$ is also bounded above by $a$. However, $(a_n)$ must be unbounded. If $(a_n)$ were bounded and monotone, it must be convergent, by the Monotone Convergence Theorem. Since $(a_n)$ is divergent, $(a_n)$ must be unbounded. So, we have a contradiction. Therefore, there cannot exist a divergent monotone sequence with a Cauchy subsequence.
    \part Let $a_n=0$ for odd $n$ and $a_n=n$ for even $n$. This is an unbounded sequence, as for every bound $M\in\mathbb{R}$, we can let $n$ be the smallest even integer greater than $M$, which would satisfy $|a_n|>M$. A subsequence of $(a_n)$ that is Cauchy would be the odd-indexed terms of $(a_n)$. Since all of the terms are 0, the limit of the subsequence is 0 and by the Cauchy Criterion, this convergent subsequence is Cauchy. So, this sequence satisfies being an unbounded sequence containing a Cauchy subsequence.
    \newline
\end{parts}

\underline{Exercise 2.6.4}
\newline
\begin{parts}
    \part To be Cauchy, $(c_n)$ must satisfy that for every $\epsilon>0$, there must exist $N\in\mathbb{N}$ such that for all $m,n\ge N$, $|c_n-c_m|<\epsilon$. Since $(a_n)$ is Cauchy, we have for every $\epsilon>0$, there exists $N_1\in\mathbb{N}$ such that for all $m,n\ge N_1$, $|a_n-a_m|<\frac{\epsilon}{2}$. Also, since $(b_n)$ is Cauchy, we have for every $\epsilon>0$, there exists $N_2\in\mathbb{N}$ such that for all $m,n\ge N_2$, $|b_n-b_m|<\frac{\epsilon}{2}$. Choose $N=\max\{N_1,N_2\}$; these conditions still hold for this choice of $N$. Then, we have $|c_n-c_m|=||a_n-b_n|-|a_m-b_m||\le|(a_n-b_n)-(a_m-b_m)|=|(a_n-a_m)-(b_n-b_m)|\le|a_n-a_m|+|b_n-b_m|<\frac{\epsilon}{2}+\frac{\epsilon}{2}=\epsilon$. Therefore, for every $\epsilon>0$, we have found a $N\in\mathbb{N}$ such that for all $m,n\ge N$, $|c_n-c_m|<\epsilon$, so $(c_n)$ is Cauchy.
    \part $(c_n)$ is not Cauchy. For any Cauchy sequence $(a_n)$ that converges to a limit that is not equal to zero, $c_n=(-1)^na_n$ does not converge. For example, let $a_n=1+\frac{1}{n}$. This is a Cauchy sequence, as it converges to 1 as $n\to\infty$. Then, $c_n=(-1)^n(1+\frac{1}{n})=(-1)^n+\frac{(-1)^n}{n}$. As $n\to\infty$, the second term is neglible, leaving $c_n\approx(-1)^n$, which does not converge to a real number. Therefore, $(c_n)$ is not convergent and not Cauchy.
    \part $(c_n)$ is not Cauchy. Let $a_n=\frac{(-1)^n}{n}$. This is a convergent sequence with a limit of 1, so it is Cauchy. Then, $c_n=[[\frac{(-1)^n}{n}]]$. For odd $n$, $-1\le a_n<0$, so $c_n=-1$. For even $n$, $0\le a_n<1$, so $c_n=0$. So, $c_n=-1$ for odd $n$ and $c_n=0$ for even $n$. This is an oscillating sequence, so $(c_n)$ does not converge, so it is not Cauchy.
    \newline
\end{parts}

\underline{Exercise 2.7.2}
\newline
\begin{parts}
    \part For all $n\in\mathbb{N}$, $0\le\frac{1}{2^n+n}\le\frac{1}{2^n}$, so we can use the Comparison Test. $\sum_{n=1}^{\infty}\frac{1}{2^n}=\sum_{n=0}^{\infty}(\frac{1}{2})^n-(\frac{1}{2})^0=\frac{1}{1-\frac{1}{2}}-1=1$. Since $\sum_{n=1}^{\infty}\frac{1}{2^n}$ is convergent, $\sum_{n=1}^{\infty}\frac{1}{2^n+n}$ is convergent by the Comparison Test.
    \part If we can show that $\sum_{n=1}^{\infty}|\frac{sin(n)}{n^2}|$ is convergent, then we will have that $\sum_{n=1}^{\infty}\frac{\sin(n)}{n^2}$ is convergent as well by the Absolute Convergence Test. $|\frac{sin(n)}{n^2}|=\frac{|sin(n)|}{n^2}<\frac{1}{n^2}<\frac{1}{n(n-1)}$ for $n>2$, so we can use the Comparison Test. Let $s_N$ represent the partial sum $\sum_{n=2}^{N}\frac{1}{n(n-1)}$. This is equal to $\sum_{n=2}^{N}\frac{1}{n-1}-\frac{1}{n}=(1-\frac{1}{2})+(\frac{1}{2}-\frac{1}{3})+\dots+(\frac{1}{N-2}-\frac{1}{N-1})+(\frac{1}{N-1}-\frac{1}{N})=1-\frac{1}{N}$. As $N\to\infty$, the limit of the sequence of partial sums goes to 1, so $\sum_{n=2}^{\infty}\frac{1}{n(n-1)}$. Then, $\sum_{n=2}^{\infty}\frac{1}{n^2}$ is convergent, and adding the term for $n=1$ keeps the series convergent. Then, by the Comparison and Absolute Convergence Tests, $\sum_{n=1}^{\infty}\frac{\sin(n)}{n^2}$ is convergent.
    \part The given series is equal to $\sum_{n=1}^{\infty}\frac{(-1)^{n+1}(n+1)}{2n}$. This is equal to $\sum_{n=1}^{\infty}\frac{(-1)^{n+1}n}{2n}+\sum_{n=1}^{\infty}\frac{(-1)^{n+1}}{2n}=\frac{1}{2}\sum_{n=1}^{\infty}(-1)^{n+1}+\frac{1}{2}\sum_{n=1}^{\infty}\frac{(-1)^{n+1}}{n}$ by the Algebraic Limit Theorem for Series. For the first sum, the sequence representing the partial sum up to the $N$th term is given by $s_N=\sum_{n=1}^{N}(-1)^{n+1}=\frac{1}{2}-\frac{1}{2}(-1)^N$. This sequence oscillates and is therefore divergent, so the sum does not converge. Therefore, the entire series does not converge.
    \part The given series can be written as $\sum_{n=1}^{\infty}\frac{1}{3n-2}+\frac{1}{3n-1}-\frac{1}{3n}$. After combining terms, we have $\sum_{n=1}^{\infty}\frac{9n^2-2}{(3n-2)(3n-1)(3n)}$. Then, $\sum_{n=1}^{\infty}\frac{9n^2-2}{(3n-2)(3n-1)(3n)}\ge\sum_{n=1}^{\infty}\frac{9n^2-2}{(3n)^3}\ge\sum_{n=1}^{\infty}\frac{8n^2}{27n^3}=\frac{8}{27}\sum_{n=1}^{\infty}\frac{1}{n}$, assuming using the terms for $n\ge2$. Since the harmonic series is divergent, the original series must be divergent by the Comparison Test.
    \part We can rewrite the given series as the sum of the positive terms minus the sum of the negative terms. Therefore, the series is equal to $\sum_{n=1}^{\infty}\frac{1}{2n-1}-\sum_{n=1}^{\infty}\frac{1}{(2n)^2}$. The first sum is greater than $\sum_{n=1}^{\infty}\frac{1}{2n}$, which is equal to $\frac{1}{2}\sum_{n=1}^{\infty}\frac{1}{n}$. Since the harmonic series is divergent, a scalar multiple of it is also divergent, so $\sum_{n=1}^{\infty}\frac{1}{2n-1}$ is divergent. To prove that the sum is divergent, the second sum must be convergent so that the difference of sums is divergent. $\sum_{n=1}^{\infty}\frac{1}{(2n)^2}=\frac{1}{4}\sum_{n=1}^{\infty}\frac{1}{n^2}$. This sum has been shown to converge in part (b) with a comparison to $\sum_{n=1}^{\infty}\frac{1}{n(n-1)}$. Therefore, the scalar multiple of the sum is convergent, so the difference of the sums is divergent. Therefore, the series is divergent.
    \newline
\end{parts}

\underline{Exercise 2.7.10}
\newline
\begin{parts}
    \part The infinite product given is equal to $\prod_{n=0}^{\infty}\frac{2^n+1}{2^n}$. This is equal to $\prod_{n=0}^{\infty}1+\frac{1}{2^n}$. From Exercise 2.4.10, products of the form $\prod_{n=1}^{\infty}1+a_n$ converge if and only if $\sum_{n=1}^{\infty}a_n$ converges. Leaving out the $n=0$ term from the product, we know $\prod_{n=1}^{\infty}1+\frac{1}{2^n}$ converges, as $\sum_{n=1}^{\infty}\frac{1}{2^n}$ converges to 1. Including the $n=0$ term just multiplies the total product by 2, so the product is convergent.
    \part The given product is equal to $\prod_{n=1}^{\infty}\frac{2n-1}{2n}$. The $n$th term satisfies $0<\frac{2n-1}{2n}<1$, so the product must converge to a limit that is at least 0. Rewriting the product, we have $\prod_{n=1}^{\infty}1-\frac{1}{2n}$. This is equal to $\prod_{n=1}^{\infty}1-a_n$, where $a_n=\frac{1}{2n}$. Using the fact that $1-x\le e^{-x}$, we can say that the product is less than or equal to $\prod_{n=1}^{\infty}e^{-a_n}$, which is equal to $e^{-\sum_{n=1}^{\infty}a_n}$. $a_n=\frac{1}{2}\frac{1}{n}$. Since the harmonic series diverges to infinity, half of it does as well, so $-\sum_{n=1}^{\infty}a_n$ goes to negative infinity. Then, raising $e$ to this value causes creates an upper bound for the product, as the product is less than or equal to this value. Since the product is greater than or equal to 0, and it is bounded above by a sequence that is tending to 0, it must be equal to 0.
    \part The product given is equal to $\prod_{n=1}^{\infty}\frac{(2n)^2}{(2n-1)(2n+1)}$. This is equal to $\prod_{n=1}^{\infty}\frac{(2n)^2}{(2n)^2-1}=\prod_{n=1}^{\infty}\frac{4n^2}{4n^2-1}=\prod_{n=1}^{\infty}1+\frac{1}{4n^2-1}=\prod_{n=1}^{\infty}1+a_n$ for $a_n=\frac{1}{4n^2-1}$. By the conclusion from Exercise 2.4.10, for the product to converge, we must prove that $\sum_{n=1}^{\infty}a_n$ converges. Since, $4n^2-1\ge n^2$ for $n\ge 1$, we know that $\frac{1}{4n^2-1}\le \frac{1}{n^2}$. Since $\sum_{n=1}^{\infty}\frac{1}{n^2}$ converges, as shown in Exercise 2.7.2.(b), the original sum must also converge. So, the product is convergent.
\end{parts}

\end{document}